\section*{الفصل \ref{ch:b3:electronic-spectroscopy} --- المطيافية الإلكترونية والفيزياء الضوئية}

\begin{solution}{exo:b3:electronic-spectroscopy:1}
(أ) ممنوع بقاعدة السبين ($\Delta S = 1$). (ب) ممنوع \omterm{thm:b3:electronic-spectroscopy:laporte}{بقاعدة لابورت}
($g \to g$)؛ ويُرى ضعيفًا عبر الاقتران \omterm{def:b3:electronic-spectroscopy:vibronic}{الاهتزازي الإلكتروني}. (ج) مسموح: $B_{1u}$ هو
تمثيل $z$ في $D_{2h}$. (د) ممنوع بالتناظر: $A_2$ ليس أيًا من $x$
و$y$ و$z$ في $C_{2v}$؛ والحزمة ضعيفة.
\end{solution}

\begin{solution}{exo:b3:electronic-spectroscopy:2}
$10^7/349 - 10^7/450 = 28653 - 22222 = \qty{6431}{cm^{-1}}$، أي \qty{0.80}{eV}.
\end{solution}

\begin{solution}{exo:b3:electronic-spectroscopy:3}
$A = 5700 \times 1 \times 1.0\times10^{-5} = 0.057$؛ والنسبة الممتصة $1 - 10^{-0.057} =
0.12$.
\end{solution}

\begin{solution}{exo:b3:electronic-spectroscopy:4}
$k_r = \Phi_F/\tau = \qty{1.5e8}{s^{-1}}$، $\tau_r = \qty{6.7}{ns}$؛ $k_{nr} = (1 - \Phi_F)/\tau
= \qty{1.0e8}{s^{-1}}$.
\end{solution}

\begin{solution}{exo:b3:electronic-spectroscopy:5}
$\int\varepsilon\,\dd\tilde\nu = 1.0645 \times 5700 \times 5000 = \num{3.03e7}$، ومنه $f \approx
4.32\times10^{-9} \times 3.03\times10^7 = 0.13$: انتقال مسموح متوسط
الشدة.
\end{solution}

\begin{solution}{exo:b3:electronic-spectroscopy:6}
عامل بواسون $\eu^{-S}S^v/v!$ أعظمي عند $v = \lfloor S\rfloor$. $S = 0.3$: الخط
0--0 نفسه. $S = 1.5$: $v' = 1$، أشد بمرة ونصف من الخط 0--0. $S = 5$: $v' = 4$
و5 بالتساوي، أشد بست وعشرين مرة من الخط 0--0.
\end{solution}

\begin{solution}{exo:b3:electronic-spectroscopy:7}
$\sum k = \qty{4.2e8}{s^{-1}}$: $\tau = \qty{2.4}{ns}$، $\Phi_F = 0.24$، $\Phi_{\mathrm{ISC}} = 0.71$
(و$\Phi_{\mathrm{IC}} = 0.05$).
\end{solution}

\begin{solution}{exo:b3:electronic-spectroscopy:8}
$\Phi = 0.546 \times 0.46 \times (0.050/0.040) \times (1.4266/1.333)^2 = 0.36$.
\end{solution}

\begin{solution}{exo:b3:electronic-spectroscopy:9}
$\tau_0/\tau = 1.429$ و1.860: خطي في $[Q]$، فالعمر نفسه
يقصر: \omterm{def:b3:electronic-spectroscopy:quencher}{إخماد ديناميكي}. $K_{SV} = 43$ L/mol (من النقطتين)، $k_q = 43/8.0
\times 10^{-9} = \qty{5.4e9}{L.mol^{-1}.s^{-1}}$.
\end{solution}

\begin{solution}{exo:b3:electronic-spectroscopy:10}
\omterm{def:b3:electronic-spectroscopy:quencher}{إخماد سكوني}: الجزيئات التي ما زالت تُصدر تعيش المدة نفسها؛ ونصفها
مرتبط، في الحالة الأساسية، في معقد غير متفلور. ومع
ثابت تشارك $K_a$، تكون النسبة الحرة $1/(1 + K_a[Q])$ و$I_0/I = 1 +
K_a[Q]$، وهي صيغة شتيرن--فولمر نفسها لكن مع $\tau$ ثابت.
\end{solution}

\begin{solution}{exo:b3:electronic-spectroscopy:11}
$\tau_r = \tau/\Phi_F = \qty{14}{ns}$؛ $\Phi_T = 1 - 0.36 = 0.64$. فالعمر المقيس
يقصّره \omterm{def:b3:electronic-spectroscopy:radiationless}{العبور بين الأنظمة} المنافس؛ أما $\tau_r$ فهو ما يعطيه الإصدار
وحده.
\end{solution}

\begin{solution}{exo:b3:electronic-spectroscopy:12}
$\frac12\langle\alpha\beta - \beta\alpha|\alpha\beta + \beta\alpha\rangle = \frac12(1 + 0 - 0 - 1) = 0$، باستعمال
$\langle\alpha|\alpha\rangle = \langle\beta|\beta\rangle = 1$، $\langle\alpha|\beta\rangle = 0$ لكل إلكترون.
\omterm{def:b3:quantum-model-systems:operator}{ومؤثر} ثنائي القطب لا يمسّ السبين، فيحتوي عزم الانتقال هذا
العامل المعدوم: لا شدة للانتقال $T_1 \leftarrow S_0$ ما لم يمزج \omterm{def:b3:many-electron-atoms:spin-orbit}{اقتران السبين--المدار}
الحالتين.
\end{solution}

\begin{solution}{pb:b3:electronic-spectroscopy:1}
\textbf{1.} $A = 5700 \times 1 \times 1.0\times10^{-4} = 0.57$.
\textbf{2.} $1 - 10^{-0.57} = 0.73$.
\textbf{3.} $\varepsilon$ من بضعة آلاف: انتقال $\pi\to\pi^*$ مسموح، متوسط
الشدة.
\textbf{4.} $1239.8/349 = \qty{3.55}{eV}$.
\textbf{5.} لا يمتص إلا في الأشعة فوق البنفسجية؛ والضوء المرئي يمر.
\textbf{6.} كي تبقى كمية الضوء الممتص متناسبة مع التركيز ولتجنب
إعادة امتصاص الضوء المصدر (مفاعيل المرشح الداخلي).
\textbf{7.} $S_0 \to S_1$ (امتصاص، ثم \omterm{def:b3:electronic-spectroscopy:radiationless}{استرخاء اهتزازي})؛ ومن $S_1$:
\omterm{def:b3:electronic-spectroscopy:luminescence}{التفلور}، \omterm{def:b3:electronic-spectroscopy:radiationless}{والتحول الداخلي}، \omterm{def:b3:electronic-spectroscopy:radiationless}{والعبور بين الأنظمة} إلى $T_1$.
\textbf{8.} $k_r = 0.546/19 \times 10^{-9} = \qty{2.87e7}{s^{-1}}$؛ $\tau_r = \qty{35}{ns}$.
\textbf{9.} $(1 - 0.546)/\tau_0 = \qty{2.39e7}{s^{-1}}$.
\textbf{10.} $28653 - 22222 = \qty{6431}{cm^{-1}}$.
\textbf{11.} ينطلق الإصدار من $S_1$ بعد استرخائها وينتهي على مستويات مثارة اهتزازيًا
من $S_0$، بعد أن يسترخي الجزيء والمذيب: فطاقة الفوتون
أقل، وهي هنا تكفي لتقع في المرئي.
\textbf{12.} $\Phi_F \times (349/450) = 0.42$: أي 42\,\% من الطاقة الممتصة.
\textbf{13.} ترتفع النقاط بمقدار 0.594 لكل \qty{0.010}{mol/L}: مستقيم يمر بالنقطة 1.
\textbf{14.} بالمربعات الصغرى على النقاط الخمس: الميل $K_{SV} = \qty{59.4}{L/mol}$،
والمقطع 1.000.
\textbf{15.} $k_q = K_{SV}/\tau_0 = \qty{3.1e9}{L.mol^{-1}.s^{-1}}$.
\textbf{16.} $\tau = \tau_0/(1 + 59.4 \times 0.040) = \qty{5.6}{ns}$.
\textbf{17.} بقياس الأعمار: في \omterm{def:b3:electronic-spectroscopy:quencher}{الإخماد الديناميكي} يقع $\tau_0/\tau$ على المستقيم نفسه
الذي يقع عليه $I_0/I$.
\textbf{18.} المياه الفوارة بالكينين لا تحتوي كلوريدًا يُخمده، وهي حمضية بما يكفي
لإبقاء الكينين مبرتنًا، وهي صورته المتفلورة.
\textbf{19.} $k_d = 8 \times 8.314 \times 298/(3 \times 0.890\times10^{-3}) =
\qty{7.4e6}{m^3.mol^{-1}.s^{-1}} = \qty{7.4e9}{L.mol^{-1}.s^{-1}}$.
\textbf{20.} $k_q$ أقل بقليل من نصف $k_d$.
\textbf{21.} نحو 42\,\%.
\textbf{22.} أبطأ: $k_d \propto 1/\eta$، و$k_q$ لا يمكن إلا أن ينخفض معه.
\textbf{23.} \textbf{$k_q/k_d = 0.42$}: يُخمد الكلوريد الكينين بسرعة تقارب نصف
سرعة اللقاءات التي يتحكم فيها الانتشار.
\end{solution}
